\documentclass{article}

% Pass options to natbib
\PassOptionsToPackage{numbers, compress}{natbib}

% Preprint option for submission/report
\usepackage[preprint]{neurips_2025}

\usepackage[utf8]{inputenc}
\usepackage[T1]{fontenc}
\usepackage{hyperref}
\usepackage{url}
\usepackage{booktabs}
\usepackage{amsfonts}
\usepackage{nicefrac}
\usepackage{microtype}
\usepackage{xcolor}
\usepackage{graphicx}
\usepackage{amsmath}

\title{\textbf{Search and Verify: Efficient Small Object Tracking in 4K Video via
Motion-Gated Native Tiling}}

\author{
  \normalsize Department of Computer Science\\
  \normalsize University of Oxford
}
\date{}

\begin{document}

\maketitle

\begin{abstract}
  Tracking small objects in high-resolution video remains a significant challenge due to
  the loss of spatial
  information during standard image resizing. In this paper, we address the task of
  tracking birds in 4K
  video sequences under realistic CPU-only compute constraints. We evaluate thirteen
  distinct strategies
  ranging from naive deep learning baselines to hybrid pipelines involving Global Motion
  Compensation (GMC),
  dynamic statistical thresholding, and temporal interpolation. Our findings demonstrate that while
  brute-force Slicing Aided Hyper Inference (SAHI) dramatically increases recall from
  0.10\% to 54.50\%, it
  incurs a heavy computational penalty. We propose Strategy 8 (Motion-Guided ROIs) and
  Strategy 12 (Temporal
  Interpolation) as optimal solutions, achieving the highest F1-score (8.29\%) and
  highest throughput (3.54
  FPS), respectively.
\end{abstract}

\section{Introduction}
Tracking small targets, such as birds in 4K video, is a fundamental problem in computer vision with
applications in conservation and wind farm safety \cite{mva2025smot4sb}. The primary
difficulty lies in the
scale; a bird may occupy only a few pixels in a $3840 \times 2160$ frame
\cite{mva2025smot4sb}. Standard
detectors like YOLO typically resize input to $640 \times 640$, causing these tiny
features to vanish. This
assignment explores methods to preserve native resolution while maintaining computational
efficiency on
CPU-only hardware \cite{clark2025vis}.

\section{Methodology}
Our methodology follows a ``Search and Verify'' paradigm. We conducted extensive
experiments involving many
permutations of strategies, confidence thresholds (0.01 to 0.20), and variations in
sharding and merging. We
define a singular true baseline for comparison, with all other methods treated as
distinct experimental strategies.

\subsection{Baseline: baseline\_base}
The true baseline for all experiments is the \textbf{baseline\_base} pipeline. It
performs standard YOLOv12n
inference on 4K frames resized directly to $640 \times 640$. This establishes the lower
bound for performance
using native deep learning techniques without specialized handling for small objects.

\subsection{Strategy 2: GMC and Dynamic Thresholding (strategy\_2.py)}
This hybrid pipeline utilizes Global Motion Compensation (GMC). It detects Shi-Tomasi
corners and tracks them
between $Frame_{t-1}$ and $Frame_t$ using Lucas-Kanade Optical Flow. By canceling camera
motion, we isolate
moving objects using an adaptive threshold $T = \mu + k\sigma$.

\subsection{Native Tiling (baseline\_w\_tiling)}
To mitigate resolution loss, we implemented a custom tiling strategy. The 4K frame is
manually divided into a
grid (e.g., 4x3) with 20\% overlap to ensure continuity across seams. Each tile is
processed at native
resolution, and detections are mapped back to global coordinates.

\subsection{SAHI Tiling (baseline\_base\_sahi)}
In this strategy, we utilize the Slicing Aided Hyper Inference (SAHI) library to handle
the sharding of the
4K frame. This serves as a comparison point for our manual native tiling implementation, testing the
efficiency of library-based patch management.

\subsection{Native Tiling + NMS (baseline\_w\_tiling\_and\_nms)}
This is an extension of the Native Tiling strategy that adds a Global Non-Maximum
Suppression (NMS) pass
after coordinate mapping. This step is critical for merging duplicate detections that
occur in the overlap
regions between tiles.

\subsection{Strategy 8: Motion-Guided ROIs and Cascaded Gates}
Strategy 8 employs GMC to identify moving blobs, generating Regions of Interest (ROIs)
that are expanded by
2.0x for context. Inference is throttled to every 5 frames, with a custom Object Tracker
maintaining state. A
full frame scan is performed every 10 frames to recover from motion-mask failures.

\subsection{Strategy 10: Motion-Gated Native Tiling}
This strategy applies native 640x640 tiling but only processes tiles flagged with
significant motion via GMC
stabilization. This focuses compute on active regions of the 4K frame.

\subsection{Strategy 10 SAHI: Motion-Gated SAHI Tiling}
This variation of Strategy 10 utilizes the SAHI library to generate the tiles that are
subsequently gated by
the motion detection logic.

\subsection{Strategy 11: Classifier before Detection}
Strategy 11 adds a ``Fail Fast'' logic. Candidate ROIs from motion detection are passed
to a lightweight
YOLOv11n-cls classifier. If the classifier does not identify a bird, the heavy detector
pass is skipped.

\subsection{Strategy 12: Interpolation with GMC}
Based on Strategy 8, this method adds \textbf{Linear Interpolation}. By running inference
every 5 frames and
skipping the middle 4, it mathematically predicts the intermediate locations of a bird if
it is detected in
both the current and previous keyframes. Once every 10 frames a full frame scan is
performed to recover from
motion-mask failures.

\subsection{Strategy 12 SAHI: Interpolation with SAHI}
This variation swaps the GMC-based ROI generation for SAHI tiling while maintaining the
linear interpolation
logic between skipped frames.

\subsection{Strategy 13: Kitchen sink with Native tiling}
Our most complex funnel, combining native tiling, GMC motion detection, and
classification. It processes
tiles individually; if motion is detected, it goes to the detector; if not, a classifier
checks for stationary birds.

\subsection{Strategy 13: Kitchen sink with Native tiling and interpolation}
This version of the "Kitchen sink" adds the temporal linear interpolation to the Native
tiling funnel to
further boost throughput.

\subsection{Strategy 13 SAHI: Kitchen sink with SAHI tiling}
The full funnel implementation using the SAHI library for initial image slicing instead
of manual native tiling.

\subsection{Strategy 13 SAHI: Kitchen sink with SAHI tiling and interpolation}
The most advanced variant, utilizing SAHI slicing, the multi-gate funnel
(motion/classification), and linear
interpolation for maximum theoretical recall and efficiency.

\section{Experiments and Results}

\subsection{Experimental Setup}
Experiments were conducted on the SMOT4SB dataset using a Google Cloud VM
(c4-standard-48). We evaluated
thresholds of 0.01, 0.05, 0.1, and 0.2 across all strategies.

\subsection{Performance Analysis}
Table \ref{results-table} summarizes the most significant strategies compared to the baseline.

\begin{table}[h]
  \caption{Performance Benchmarks on SMOT4SB Dataset}
  \label{results-table}
  \centering
  \small
  \begin{tabular}{lllll}
    \toprule
    \textbf{Strategy} & \textbf{Avg FPS} & \textbf{Precision} & \textbf{Recall} &
    \textbf{F1-Score} \\
    \midrule
    Baseline Base (0.01) & 0.69 & 3.41\% & 0.10\% & 0.19\% \\
    Baseline Base + SAHI & 0.44 & 2.39\% & 54.50\% & 4.58\% \\
    Strategy 2 + SAHI & 0.03 & 2.15\% & \textbf{62.30\%} & 4.15\% \\
    Strategy 8 + SAHI & 1.82 & 4.65\% & 38.66\% & \textbf{8.29\%} \\
    Strategy 12 (0.01) & \textbf{3.54} & 7.85\% & 0.87\% & 1.56\% \\
    Strategy 12 + SAHI & 2.25 & 2.53\% & 43.89\% & 4.78\% \\
    \bottomrule
  \end{tabular}
\end{table}

\subsection{Discussion}
\textbf{The Slicing Paradox:} SAHI remains the most effective method for increasing
recall (+54,400\% over
baseline), but its 0.44 FPS throughput is unsuitable for real-time use. \textbf{Temporal
Efficiency:}
Strategy 12 achieved \textbf{3.54 FPS} using linear interpolation.

\section{Conclusion}
Our results indicate that \textbf{Strategy 8 with SAHI} provides the best balance of
detection quality
(8.29\% F1), while \textbf{Strategy 12} offers the highest efficiency for CPU-based tracking.

\bibliographystyle{unsrtnat}
\bibliography{refs}

\end{document}